\section{Sheafification of abelian presheaves}
\label{section-sheafify-abelian-presheaves}

\noindent
The following strange looking lemma is likely unnecessary, but
very convenient to deal with sheafification of presheaves
of algebraic structures.

\begin{lemma}
\label{lemma-diagram-fibre-product}
Let $X$ be a topological space. Let $\mathcal{F}$ be
a presheaf of sets on $X$. Let $U \subset X$ be open.
There is a canonical fibre product diagram
$$
\xymatrix{
\mathcal{F}^\#(U) \ar[d] \ar[r] &
\Pi(\mathcal{F})(U) \ar[d] \\
\prod_{x \in U} \mathcal{F}_x
\ar[r] &
\prod_{x \in U} \Pi(\mathcal{F})_x
}
$$
where the maps are the following:
\begin{enumerate}
\item The left vertical map has components
$\mathcal{F}^\#(U) \to \mathcal{F}^\#_x = \mathcal{F}_x$
where the equality is Lemma \ref{lemma-stalk-sheafification}.
\item The top horizontal map comes from the
map of presheaves $\mathcal{F} \to \Pi(\mathcal{F})$ described
in Section \ref{section-sheafification}.
\item The right vertical map has obvious component
maps $\Pi(\mathcal{F})(U) \to \Pi(\mathcal{F})_x$.
\item The bottom horizontal map has components
$\mathcal{F}_x \to \Pi(\mathcal{F})_x$
which come from the map of presheaves
$\mathcal{F} \to \Pi(\mathcal{F})$ described
in Section \ref{section-sheafification}.
\end{enumerate}
\end{lemma}

\begin{proof}
It is clear that the diagram commutes. We have to show
it is a fibre product diagram. The bottom horizontal arrow
is injective since all the maps $\mathcal{F}_x \to \Pi(\mathcal{F})_x$
are injective (see beginning proof of
Lemma \ref{lemma-stalk-sheafification}).
A section $s \in \Pi(\mathcal{F})(U)$ is in $\mathcal{F}^\#$ if and
only if $(*)$ holds. But $(*)$ says that around every point
the section $s$ comes from a section of $\mathcal{F}$. By definition
of the stalk functors, this is equivalent to saying that
the value of $s$ in every stalk $\Pi(\mathcal{F})_x$ comes
from an element of the stalk $\mathcal{F}_x$. Hence the lemma.
\end{proof}


\begin{lemma}
\label{lemma-sheafify-abelian-presheaf}
Let $X$ be a topological space.
Let $\mathcal{F}$ be an abelian presheaf on $X$.
Then there exists a unique structure of
abelian sheaf on $\mathcal{F}^\#$ such that
$\mathcal{F} \to \mathcal{F}^\#$ is a morphism
of abelian presheaves. Moreover, the following adjointness
property holds
$$
\Mor_{\textit{PAb}(X)}(\mathcal{F}, i(\mathcal{G}))
=
\Mor_{\textit{Ab}(X)}(\mathcal{F}^\#, \mathcal{G}).
$$
\end{lemma}

\begin{proof}
Recall the sheaf of sets $\Pi(\mathcal{F})$ defined in
Section \ref{section-sheafification}. All the stalks
$\mathcal{F}_x$ are abelian groups, see
Lemma \ref{lemma-stalk-abelian-presheaf}.
Hence $\Pi(\mathcal{F})$ is a sheaf of abelian groups by
Example \ref{example-sheaf-product-pointwise-algebraic-structure}.
Also, it is clear that the map $\mathcal{F} \to \Pi(\mathcal{F})$
is a morphism of abelian presheaves. If we show that
condition $(*)$ of Section \ref{section-sheafification} defines a subgroup
of $\Pi(\mathcal{F})(U)$ for all open subsets $U \subset X$,
then $\mathcal{F}^\#$ canonically inherits the structure of abelian sheaf.
This is quite easy to do by hand, and we leave it to the
reader to find a good simple argument. The argument we use here,
which generalizes to presheaves of algebraic structures is the following:
Lemma \ref{lemma-diagram-fibre-product} show that
$\mathcal{F}^\#(U)$ is the fibre product of a diagram
of abelian groups. Thus $\mathcal{F}^\#$ is an abelian
subgroup as desired.

\medskip\noindent
Note that at this point $\mathcal{F}^\#_x$ is an abelian
group by Lemma \ref{lemma-stalk-abelian-presheaf}
and that $\mathcal{F}_x \to \mathcal{F}^\#_x$ is a
bijection (Lemma \ref{lemma-stalk-sheafification})
and a homomorphism of abelian groups. Hence
$\mathcal{F}_x \to \mathcal{F}^\#_x$ is an isomorphism
of abelian groups. This will be used below without further mention.

\medskip\noindent
To prove the adjointness property we use the adjointness
property of sheafification of presheaves of sets. For example
if $\psi : \mathcal{F} \to i(\mathcal{G})$ is morphism of presheaves
then we obtain a morphism of sheaves
$\psi' : \mathcal{F}^\# \to \mathcal{G}$. What we have to do is to check
that this is a morphism of abelian sheaves. We may do this
for example by noting that it is true on stalks,
by Lemma \ref{lemma-stalk-sheafification}, and then using
Lemma \ref{lemma-check-homomorphism-stalks} above.
\end{proof}
